% 本文件由 LaTeX 排版 Agent 根据有效 translation.json 自动生成。
% author=Clement of Rome; work_id=0803; title=论童贞两书

\chapter{论童贞两书}

% structure: p0001 source=h1 target=omitted reason=page-title-already-rendered-by-parent

% structure: p0002 source=h2 target=section reason=relative-heading-rank; base=h2

\section{第一书}

% structure: p0003 source=h3 target=subsection reason=relative-heading-rank; base=h2

\subsection{第一章 致候}

致所有爱惜并珍视他们\Emph{在基督内的}生命、借着天主父，且怀着永生的希望服从天主的真理的人；致那些在爱德中向弟兄和近人怀有感情的人；致蒙福的\Emph{弟兄}童贞者，他们奉献自己为保守童贞\BibleQuote{为天国；} \BibleRef{玛 19:12}以及至圣的\Emph{姊妹}童贞者：愿天主的平安与你们同在。\par

% structure: p0005 source=h3 target=subsection reason=relative-heading-rank; base=h2

\subsection{第二章 真童贞须有全德}

凡男女童贞者中，凡真正决心为天国而保守童贞的，每一位都必须凡事堪当天国。因为天国不是靠口才或名声、地位或出身、美貌或力量、或长寿获得的；而是靠信德的能力，当人显出信德的作为时获得的。因为凡真正义的人，他的作为证实他的信德，证明他确实是信徒，具有伟大的信德、完美的信德、在于天主的信德、在善行中发光的信德，以使万有的父藉着基督受光荣。如今，那些为天主而真正童贞的人听从那说过\Quote{愿正义和信德不离开你；将它们系在你的颈上，你必为自己寻得恩宠，在天主和人面前筹谋善事。}因此，\Quote{义人的道路好像黎明的光，越来越亮，直到日午。}因为他们的光芒甚至现在已藉善行光照整个受造界，作为那些真正\BibleQuote{世界的光，} \BibleRef{玛 5:14}光照\Quote{那些坐在黑暗中的人，}使他们能起来并从黑暗出来，藉着敬畏天主的善行之光，\Quote{好叫他们看见我们的善行，光荣我们在天之父。}因为天主的人必须在一切言语和行为上完美，在生活上以一切模范和有序的品行装饰，并在正义中行一切事，作为天主的人。\par

% structure: p0007 source=h3 target=subsection reason=relative-heading-rank; base=h2

\subsection{第三章 真童贞者藉克己证明自己，正如真信徒藉善行证明一样}

因为童贞者是信徒和将要相信的人的美丽典范。但仅有名称而无作为，并不能引入天国；但如果人真是信徒，这样的人可以得救。因为如果一个人只在名义上被称为信徒，而在行为上却不是，他就不可能是一个信徒。因此，\Quote{不要让任何人}因此\Quote{以虚空的言语迷惑你们。} \BibleRef{弗 5:6}因为仅仅被称为童贞者，若缺乏与童贞相称的卓越、合宜的作为，他不可能得救。因为我们的主称那样的童贞为\Quote{愚拙的，} \BibleRef{玛 25:2}因为它既没有油也没有光，就被遗留在天国之外，被关在新郎的喜乐之外，并被视为与祂的仇敌一道。因为这样的人们\BibleQuote{只有敬虔的外貌，却否认敬虔的能力。} \BibleRef{弟后 3:5}因为他们\Quote{自以为什么，其实什么都不是，乃是自欺了。但各人当不断省察自己的作为，} \BibleRef{迦 6:3}并认识自己；因为那宣认童贞和圣洁却\Quote{否认其能力}的人，献上的是空洞的敬拜。因为这样的童贞是不洁的，被一切善行所否认。因为\Quote{凡树由其果实可知。} \Quote{你要思想我所说的，主必赐你领悟。} \BibleRef{弟后 2:7}因为凡在天主面前立志保持圣洁的人，必须以天主的一切圣洁能力束腰。并且若他怀着真实的敬畏而钉死自己的身体，便因敬畏天主而免除\Emph{其中}\Emph{圣经}所说的那句话：\BibleQuote{你们要生育繁殖，} \BibleRef{创 1:28}并\Emph{回避}这世界的一切炫耀、忧虑、感官享受和迷惑，以及它的宴乐和醉酒，一切奢侈和安逸；并从这世界的整个生活及其网罗、陷阱和阻碍中退出；当你行走在地上时，要切望你的工作与事业在天上。\par

% structure: p0009 source=h3 target=subsection reason=relative-heading-rank; base=h2

\subsection{第四章 论克己的继续；真童贞者的目标与赏报}

因为那为自己贪求这些\Emph{如此}伟大卓越之事的人，便因此抽身并断绝与整个世界的联系，以便去\Emph{并}度一种神圣属天的生活，如圣洁的天使，在纯洁圣善的作为中，\BibleQuote{在天主圣神的圣洁中，} \BibleRef{得后 2:13}并为天国的缘故，藉着耶稣基督事奉全能的天主。为此，他断绝一切肉体的欲望。不仅他免除这条\Emph{命令}：\BibleQuote{你们要生育繁殖，} \BibleRef{创 1:28}而且他渴望那\BibleQuote{所应许的、在天上所储藏的希望，} \BibleRef{哥 1:5}由天主亲口宣布（祂不撒谎）：\BibleQuote{胜过儿女；} \BibleRef{依 56:4}并且祂要赐给童贞者在天主家中一个显赫的地位，这乃是\Emph{某种}\BibleQuote{胜过儿女的事，}并且比那些在圣洁中度过婚姻生活者的\Emph{地位}更好，他们的\BibleQuote{床没有玷污。} \BibleRef{希 13:4}因为天主将把天国赐给童贞者，如圣洁的天使一样，因这伟大而高尚的宣认。\par

% structure: p0011 source=h3 target=subsection reason=relative-heading-rank; base=h2

\subsection{第五章 童贞的苦辛与仇敌}

那麼，你願意作貞女嗎？你知道真正的貞潔——那時時刻刻常立於天主面前、不從祂的服務退縮、\Quote{掛慮如何以聖潔的身體和精神取悅其主}——有何等艱辛與令人厭煩嗎？\BibleRef{格前 7:34} 你知道貞潔有何等大的榮耀，並為此而\Emph{決意}修練嗎？你真的知道並明白你所渴望的是什麼嗎？你熟悉聖潔貞潔的崇高任務嗎？你知道如何像男子一般\Quote{按規矩}進入這場競賽並\Quote{較量}\BibleRef{弟後 2:5}，好藉聖神的力量，為自己選擇此道，使你戴上光明的冠冕，並被人\Emph{凱旋}引領穿越\Quote{天上的耶路撒冷}嗎？\BibleRef{迦 4:26} 因此，若你渴求這一切，就制伏身體；制伏肉體的情慾；以天主的神制伏世界；制伏這些虛幻的、短暫的、會消逝、衰敗、朽壞、終結的事物；制伏巨龍；\BibleRef{默 12:7} 制伏獅子；\BibleRef{伯前 5:8} 制伏蛇；\BibleRef{格後 11:3} 制伏撒殫——藉著耶穌基督，祂藉著祂言語的聽聞和神聖的聖體聖事堅固你。\Quote{背起你的十字架跟隨}\BibleRef{瑪 16:24}那位潔淨你的主耶穌基督。努力直奔向前，勇敢無畏，不膽怯，卻懷著勇氣，依賴你主的應許，你必獲得那\Quote{在上召叫}\BibleRef{斐 3:14}的勝利冠冕，藉著耶穌基督。因為凡在信德中行走完全、無所畏懼的人，確實領受貞潔的冠冕，它的勞苦偉大，賞報也偉大。你明白並知道聖潔是何等可敬的事嗎？你知道貞潔的榮耀是何等偉大、崇高、卓越嗎？\par

% structure: p0013 source=h3 target=subsection reason=relative-heading-rank; base=h2

\subsection{第六章　貞潔的神性}

一位神聖貞女的子宮孕育了我們的主耶穌基督、天主子；我們的主所穿戴的身體，並在其中於此世作戰的身體，祂是從一位神聖貞女取得的。因此，要明白貞潔的偉大與尊嚴。你願意作基督徒嗎？在一切事上效法基督。若翰，那在我們主之前來臨的使者，他\Quote{在婦女所生者中沒有更大的}\BibleRef{瑪 11:11}，我們主的聖使，曾是一位貞潔者。因此，效法我們主的使者，在一切事上作他的跟隨者。又那位\Quote{斜靠在我們主懷裡、祂所深愛的}\BibleRef{若 21:20}若翰——他同樣是一位聖潔的人。因為我們主愛他，並非沒有理由。保祿、巴爾納伯、弟茂德，以及所有其餘\Quote{名字寫在生命冊上}\BibleRef{斐 4:3}的人——我說，這些人都珍視並喜愛聖潔，在競賽中奔跑，無可指責地完成了他們的賽程，作為基督的效法者，作為永生天主的子女。此外，我們也發現厄里亞、厄里叟，以及許多其他聖潔的人，度過了聖潔無瑕的生活。因此，若你願意效法他們，就盡全力效法他們。因為經上說：\Quote{你們要尊敬在你們中間的長老，並觀看他們的生活方式和品行，效法他們的信德。}\BibleRef{希 13:7} 又說：\Quote{我的弟兄們，你們要效法我，如同我\Emph{效法}基督一樣。}\BibleRef{格前 11:1}\par

% structure: p0015 source=h3 target=subsection reason=relative-heading-rank; base=h2

\subsection{第七章　真正的貞女}

因此，凡效法基督的人，就誠懇地效法祂。因為那些\Quote{穿上基督}\BibleRef{羅 13:14}的人，真實地在他們的思想、整個生命和一切行為中表達祂的模樣：在言語、行為、忍耐、剛毅、知識、貞潔、含忍、純潔的心、信德、望德，以及對天主圓滿而完美的愛德中。所以，除非一位貞女在一切事上都像基督，並像\Quote{屬於基督的人}\BibleRef{迦 5:24}，否則不能得救。因為凡在天主內的貞女，她的身體和靈魂都是聖潔的，並恆心於事奉她的主，不偏離任何方向，卻常在純潔和聖潔中，在天主的神裡等候祂，\Quote{掛慮如何取悅她的主}\BibleRef{格前 7:32}，\Emph{藉著}純潔無瑕地生活，並渴望在一切事上討祂喜悅。這樣的貞女不離開我們的主，卻在精神上\Emph{常}與她的主同在：正如經上所記：\Quote{你們要聖潔，因為我是聖潔的，上主說。}\par

% structure: p0017 source=h3 target=subsection reason=relative-heading-rank; base=h2

\subsection{第八章　貞女藉著摒除一切肉體的情感，成為天主的效法者}

因为，如果有人只在名义上被称为圣人，他并不是圣人；他必须在一切事上成为圣人：在身体和灵魂上。而那些守贞的人时刻因肖似天主及其基督而喜乐，并且是他们的效法者。因为在这些如此的人身上，没有\Quote{肉性的思念。}在那些真正相信的人，以及\Quote{基督的圣神居住在他们内}\BibleRef{罗 8:9}的人身上——在他们身上不可能有\Quote{肉性的思念}，即：奸淫、不洁、放荡；崇拜偶像、行邪术；仇恨、嫉妒、竞争、忿怒、争论、分裂、恶意；醉酒、狂欢；戏谑、愚妄的谈吐、喧闹的大笑；毁谤、暗讽；苦涩、暴怒；喧嚷、辱骂、言语无礼；恶毒、捏造恶事、谎言；多言、絮叨；恐吓、切齿、轻易指控、冲突、轻视、殴打；\Emph{正义}的扭曲、\Emph{判断}的松弛；高傲、自大、张扬、浮夸、\Emph{夸耀}家世、美貌、地位、财富、血肉的臂膀；好争吵、不义、好胜；仇恨、忿怒、嫉妒、背信、报复；荒淫、贪饕、\Quote{贪婪（即崇拜偶像），}\BibleRef{哥 3:5}\Quote{贪爱钱财（万恶之根）；}\BibleRef{弟前 6:10}爱炫耀、虚荣、好权、自以为是、骄傲（被称为死亡，且是\Quote{天主所敌对的}）。凡有这些以及类似之事的人——这样的人都是属于肉性。因为\Quote{由肉性生的就是肉性；那出于地上的，讲论的是地上的事，}\BibleRef{若 3:6}他的思念也属于地上。而且\Quote{肉性的思念是与天主为敌的。因为它不服从天主的法律；不能\Emph{如此}，}\BibleRef{罗 8:7}因为它属于肉性，\Quote{在肉性内没有善，}\BibleRef{罗 7:18}因为天主的圣神不在它内。为此，圣经合理地对这样的世代说：\Quote{我的圣神将不再永久居住在人间，因为他们是肉性。}\Quote{所以，凡没有天主圣神在他内的，就不属于祂：}正如经上所载：\Quote{天主的圣神离开了撒乌耳，便有来自天主的一个恶神扰乱他。}\BibleRef{撒上 16:14}\par

% structure: p0019 source=h3 target=subsection reason=relative-heading-rank; base=h2

\subsection{第九章 继续论克己；献身于天主者的尊严。}

凡天主圣神在他内的，都符合天主圣神的旨意；并且，因为他符合圣神，所以他克治身体的作为，为天主而生活，\Quote{制伏并驯服身体，使它服从；这样，在向他人宣讲时}，他可以成为信徒美好的榜样和模范，并在一生中从事与圣神相称的善工，使他\Quote{不至于被弃绝}\BibleRef{格前 9:27}，而是天主和众人眼前蒙悦纳。因为\Quote{属天主的人}\BibleRef{弟前 6:11}，\Emph{我敢说}，他身上没有肉性的思念；尤其在守贞的\Emph{男女}中；但他们所有的果实都是\Quote{圣神的果实}\BibleRef{迦 5:22}和生命的果实，他们确实是天主的城，是天主安居和居住的房屋和圣殿，祂在其中行走，如同在圣洁的天上之城。因为\Quote{你们在这世界上显为光明，因为你们持守生命的话语}\BibleRef{斐 2:15}，这样你们实在是我们的主耶稣基督内善仆的赞美、夸耀、喜乐的冠冕和喜悦。因为凡看见你们的，必将承认\Quote{你们是上主所祝福的苗裔}\BibleRef{依 61:9}；确实是尊贵圣洁的苗裔，\Quote{是王家的司祭、圣洁的邦国、属于继承的子民}\BibleRef{伯前 2:9}，是天主恩许的继承人；\Emph{即}那些不腐朽、不凋谢的事物；\Quote{即是眼所未见、耳所未闻、人心所未想到的，就是天主为爱祂并遵守祂诫命的人所预备的。}\BibleRef{格前 2:9}\par

% structure: p0021 source=h3 target=subsection reason=relative-heading-rank; base=h2

\subsection{第十章 谴责与少女的危险和可耻交往}

现在，我的弟兄们，我们深信你们念念不忘那些为你们得救所必需的事。但我们这样说，是因为听到了关于无耻之人的恶言和报告，他们以敬畏天主为借口，与少女同住，\Emph{如此}，将自己暴露于危险之中，与她们独自在道路和僻静之处同行——这条路充满危险，且满是绊脚石、圈套和陷阱；基督徒和敬畏天主的人无论如何都不应如此行事。另一些人，在宴席上与她们一同吃喝，\Emph{放纵自己}于放荡的行为和许多不洁——这些事在信徒中，尤其在那些为自己选择了\Emph{圣洁生活}的人中，本不应有。还有一些人，聚集一处，为的是虚浮琐碎的交谈和嬉笑，以及彼此说坏话；他们搜集彼此的故事，并且无所事事：我们甚至不许你们与这样的人一同吃饭。还有一些人，假借探访、\Emph{为他们}读经或驱魔之名，流连于守贞弟兄或姊妹的家中。既然他们游手好闲、不工作，就探询那些不该过问的事，并藉花言巧语利用基督的名做生意。\Emph{这些就是}那神圣宗徒所远离的人，因为他们恶\Emph{行}众多；正如经上所载：\Quote{懒汉手中有荆棘发芽；}\BibleRef{箴 26:9}以及\Quote{懒人的道路布满荆棘。}\par

% structure: p0023 source=h3 target=subsection reason=relative-heading-rank; base=h2

\subsection{第十一章 懒惰的危害；警告勿空想成为教师；关于教导和运用神恩的劝告。}

凡不工作、反倒猎奇寻谈、自认为此事有益且正确之人，其行径正是如此。因为这样的人就像那些懒惰、饶舌的寡妇，\BibleQuote{她们挨家游荡}\BibleRef{弟前 5:13}，喋喋不休，搜索闲谈，夸大其词地挨家传递，毫无敬畏天主之心。此外，他们厚颜无耻，以教导为借口，提出各种学说。但愿他们传授真理的学说！但令人\Emph{如此}不安的，正是他们不明白自己在说什么，断言\Emph{不真实}的事：因为他们想当教师，炫耀自己口才；他们以基督的名义贩卖不义——这是天主的仆人不\Emph{应为}之事。他们不听从圣经所说：\Quote{我的弟兄们，不要多人作教师，也不要人人作先知。}因为\BibleQuote{在言语上没有过失的，便是完人，能够约束并降服全身。}\BibleRef{雅 3:2}并且，\BibleQuote{若有人讲话，当讲天主的话。}\BibleRef{伯前 4:11}还有，\BibleQuote{你若明白，就当回答兄弟；若不，就当用手掩口。}\BibleRef{德 5:14}因为，\BibleQuote{有时\Emph{当}静默，有时当言语。}\BibleRef{训 3:7}又说：\BibleQuote{人说话合其时，是荣耀之事。}\BibleRef{箴 25:11}又说：\Quote{你们的言语要常常和蔼可亲。因为人必须知道如何按季节回答每个人。}因为\Quote{凡随口乱说之人，制造纷争；多言多语，增加烦恼；嘴唇急躁，陷于邪恶。因为舌头不驯引发怒气；但完人谨守舌头，爱护自己的生命。}因为这些人就是\BibleQuote{藉花言巧语和悦人口舌迷惑朴实人心，表面祝福，实则引人走入歧途。}\BibleRef{罗 16:17}因此，我们当畏惧那等待教师的审判。因为那些\BibleQuote{教导别人却自己不行}\BibleRef{玛 23:3}的教师，以及那些假借基督之名、取\Emph{于己}身，却说：我们教导真理，却\Emph{又}闲游浪荡，自高自大，以血肉的心思夸耀的人，必受严厉审判。\BibleRef{哥 2:18}此外，他们就像\BibleQuote{瞎子领瞎子，两人都掉进坑里。}\BibleRef{玛 15:14}他们必受审判，因为他们的饶舌和轻浮的教训教导属世的智慧，以及\Quote{花言巧语之人的轻浮谬误}，\Quote{依照那掌管空中区域之首领的意志，以及那运行在悖逆之人身上的灵，按照这世界的训练，而非基督的学说。}但若你领受了\BibleQuote{知识的言语，或教导的言语，或预言的言语}\BibleRef{格前 12:8}，愿天主受赞美，\BibleQuote{祂毫不吝惜地帮助每个人——那位赐予每个人、且不责备\Emph{他}的天主。}\BibleRef{雅 1:5}因此，以你从我们主所领受的恩赐，服事你的属灵弟兄——那些知道你所言是我们主的话语的先知；并宣告你在教会中为基督内弟兄的建树所领受的恩赐（因为那些帮助天主之人的事是美好而卓越的），倘若这些恩赐确实与你同在。\par

% structure: p0025 source=h3 target=subsection reason=relative-heading-rank; base=h2

\subsection{第十二章. 访客、驱魔、如何协助病人以及凡事无过地行走的规则}

此外，这也是合宜而有益的：人当\Quote{探望孤儿寡妇}，\BibleRef{雅 1:27}尤其那些子女众多的贫苦人。这些事毫无疑问，是天主的仆人应做的，对他们也是合宜而适当的。再者，对于在基督内的弟兄，这也是适当、正确而合宜的：他们应当探望那些被恶神困扰的人，明智地为他们祈祷并驱逐，献上天主悦纳的祈祷；不是以大量优美、精心准备和编排的言辞，为在人前显得有口才和记忆力好。这样的人是\Quote{如同鸣的锣、响的钹一样}；\BibleRef{格前 13:1}他们对所驱逐的人毫无帮助；反而以可怕的言语说话，惊吓众人，却不凭真实信德行事，依照我主教导，祂曾说：\Quote{这一类非用禁食和祈祷不能出去}，应恒心诚意地祈祷。他们当圣洁地向天主恳求，以喜乐、谨慎和纯洁，不怀仇恨和恶意。如此，我们当接近患病的弟兄或姊妹，以正确的方式探望他们，无欺诈、无贪欲、无喧哗、无多言、无违逆敬畏天主之心的行为、无傲慢，却要有基督的温柔和谦卑精神。因此，让他们以禁食和祈祷进行驱逐，不是以优雅、精心编排的学问辞藻，而是作为从天主领受治愈之恩的人，满怀信心，为光荣天主。藉你们的禁食、祈祷和不懈的警醒，连同你们的其他善工，藉圣神的能力，克制肉身的作为。这样行事的人\Quote{是天主圣神的宫殿}。\BibleRef{格前 6:19}让他驱逐魔鬼，天主必帮助他。因为帮助患病的人是好的。我主曾说：\Quote{驱逐魔鬼}，同时\Emph{命令}了许多其他的治愈事工；并且，\Quote{你们白白得来的，也要白白施舍}。\BibleRef{玛 10:8}这样的人，天主为他们预备了优厚的赏报，因为他们用主所赐的恩赐服事弟兄。这对天主的仆人也是合宜而有助的，因为他们遵从我主的命令，祂曾说：\Quote{我患病，你们探望了我，等等}。这也是合宜、正确而公正的：我们应当为天主的缘故，以一切合宜的态度和纯洁的品行探望邻人；正如宗徒所说：\Quote{有谁软弱，我不软弱呢？有谁跌倒，我不着急呢？}\BibleRef{格后 11:29}但这一切话都是论及人当如何爱邻人的爱德。我们要在这些事上殷勤，不使人跌倒，不徇私情或羞辱他人，却要爱穷人如同天主的仆人，尤其是探望他们。因为在天主和人面前，这是合宜的：我们当记念穷人，作爱弟兄和爱外乡的人，为天主的缘故，也为那些信从天主的人的缘故，正如我们从法律、先知以及我们的主耶稣基督所学到的，关于爱弟兄和爱外乡之道：因为你们知道那些关于爱弟兄和爱外乡的话；这些话有力地临到所有实行它们的人。\par

% structure: p0027 source=h3 target=subsection reason=relative-heading-rank; base=h2

\subsection{第十三章 司铎应为何等人，不应为何等人。}

亲爱的弟兄们！人应在对独一天主的信德上建立和坚固弟兄，这也是明显而众所周知的。再者，这也是合宜的：人不应嫉妒邻人。并且，再次，所有从事主的工作的人，都应在敬畏天主中从事主的工作，这是适当而合宜的。他们应当如此行事。\Quote{庄稼多，工人少}，这也是众所周知和明显的。因此，让我们\Quote{求庄稼的主人}派遣工人去收祂的庄稼；\BibleRef{玛 9:37}这样的工人应\Quote{正确无误地分施真理之言}；工人\Quote{不以福音为耻}；忠信的工人；工人应是\Quote{世界的光}；\BibleRef{玛 5:14}工人\Quote{不是为了那必朽坏的食物劳作，而是为了那存留到永生的食粮}；\BibleRef{若 6:27}工人应如同宗徒一样；工人应效法圣父、圣子和圣神；关心人的救恩；不是\Quote{佣工}；\BibleRef{若 10:12}不是那些以敬畏天主和正义为可获利的工人；不是那些\Quote{服事自己的肚腹}的工人；不是那些\Quote{以花言巧语诱惑老实人的心}的工人；\BibleRef{罗 16:18}不是那些模仿光明之子，而自己不是光却是黑暗的工人——\Quote{他们的结局是丧亡}；\BibleRef{斐 3:9}不是那些行不义、邪恶和欺诈的工人；不是\Quote{诡诈的工人}；\BibleRef{格后 11:13}不是\Quote{酗酒}和\Quote{不忠信}的工人；也不是贩卖基督的工人；不是迷惑人的；不是\Quote{贪爱钱财的；不怀恶意}的工人。\par

因此，让我们注视并效法那些在主内善行度日的忠信之人，正如与我们的召叫和使命相称而合宜的。这样，我们当在天主面前，以正义和正直，毫无瑕疵地事奉，\Quote{专心从事在天主面前以及在众人面前皆为美好和合宜的事}。\BibleRef{罗 12:17}因为天主在我们身上因一切事受光荣，这是合宜的。\par

% structure: p0030 source=h2 target=section reason=relative-heading-rank; base=h2

\section{第二封书信}

% structure: p0031 source=h3 target=subsection reason=relative-heading-rank; base=h2

\subsection{第一章 他描述自己与异性交往的谨慎，并述说在旅途中，当只有弟兄的地方，他如何行事。}

此外，我的弟兄们，我还愿意你们知道我们在基督内的品行，以及我们所有弟兄在\Emph{各种}所处之处的品行。若你们认可，你们也照样在主内行事。至于我们，若天主助佑，我们这样行事：不与童贞女同住，也不与她们有任何共同之处；不与童贞女同食共饮；在童贞女睡觉之处，我们也不睡觉；妇女也不为我们洗脚或傅油；我们绝不睡在未婚或已发愿的童贞女睡觉之处；即使她\Emph{独自}一人，我们也不在那里过夜。此外，若\Emph{休息}的时间赶上我们，无论是在乡间、村庄、城镇、小村落，或无论我们在何处，若在那里发现弟兄们，我们就投靠一位弟兄，在那里召集所有弟兄，向他们谈论鼓励和劝勉的话。我们中间有讲道恩赐的人，将说出热诚、严肃、纯洁、敬畏天主的话语，劝勉他们在一切事上中悦天主，在善工上丰盛而前进，并\BibleQuote{在一切事上无忧无虑}，\BibleRef{斐 4:6} 正如天主子民所适宜合理的。\par

% structure: p0033 source=h3 target=subsection reason=relative-heading-rank; base=h2

\subsection{第二章 他在有男女基督徒之处的行为}

此外，若我们远离家园和邻舍，白日渐逝，黄昏降临，弟兄们因手足之爱和好客之情，勉强我们与他们同住，好使我们与他们一同守望，聆听天主圣言并\Emph{实行}它，从主的话语得喂养，牢记它们；他们为我们摆上饼和水以及天主所供给的；我们愿意并同意与他们过夜。若那里有一位圣洁的人，我们就投靠他并留宿，那同一位弟兄将提供并预备我们所需要的一切；他亲自服侍我们，亲自为我们洗脚并傅以香膏，亲自为我们铺床，使我们能信赖天主安眠。这一切事都由那位在我们所在之处的奉献弟兄亲自去做。他将亲自服侍弟兄们，而同一处的每位弟兄都将与他一同提供一切对弟兄们必需的服侍。但那时不得有任何女性与我们同在，无论是年轻童贞女或已婚妇女；也不得有年迈者，或已发愿者；甚至婢女，无论是基督徒或异教徒；只可有男人与男人。再者，若我们觉得有必要为妇女们站立祈祷，并宣讲劝勉和建树的话，我们\Emph{一起}召集弟兄们和所有圣洁的姊妹及童贞女，\Emph{也}同样召集那里的所有其他妇女，以全然的谦逊和得体的举止\Emph{邀请她们}前来饱飨真理。我们中间有讲道恩赐的人向她们说话，用天主赐给我们的那些话语劝勉她们。然后我们祈祷，互相问候，男人问候男人。但妇女和童贞女要用衣服裹住双手；我们也要谨慎且全然纯洁，目光向上，用衣服裹住右手；然后她们前来，在我们裹着衣服的右手上给予问候。然后我们到天主容许我们去的地方。\par

% structure: p0035 source=h3 target=subsection reason=relative-heading-rank; base=h2

\subsection{第三章 在仅有已婚基督徒之处独身弟兄的行为规则}

若我们再次偶然来到一处没有奉献弟兄、但所有人都已婚的地方，那里所有的人都会接待到来的弟兄，服侍他，殷勤而善意地关心他的一切需要。那位弟兄将以适当的方式接受他们的服侍。那位弟兄会对那里的已婚者说：我们圣洁的人不与妇女同食共饮，也不由妇女或童贞女服侍，妇女不为我们洗脚或傅油，也不为我们铺床，我们不在妇女睡觉之处睡觉，使我们在一切事上无可指责，免得有人因我们而跌倒或绊倒。当我们遵守这一切时，\BibleQuote{我们便对众人无可指摘}。\BibleRef{格后 6:3} 因此，作为\BibleQuote{知道敬畏主，我们劝导人，而在天主面前却是显明的}。\BibleRef{格后 5:11}\par

% structure: p0037 source=h3 target=subsection reason=relative-heading-rank; base=h2

\subsection{第四章 在仅有妇女之处圣洁之人的行为}

但若我们偶然来到一处没有基督徒男子、所有信徒都是妇女和童贞女的地方，她们坚持要我们在那里过夜，我们就将她们全部召集到一个合适的地方，询问她们的情况；根据我们从她们那里了解到的，以及我们所见她们的心境，我们以适宜的方式对她们说话，如同敬畏天主的人。当她们都聚集并\Emph{一同}前来，我们见她们平安，就向她们宣讲敬畏天主的劝勉话语，以纯洁和简短而有力的敬畏天主之言向她们诵读圣经。我们一切所为都为她们的建树。至于已婚者，我们以适合她们的方式在主内对她们说话。此外，若白日渐逝，黄昏临近，我们挑选一位年长的、在她们中间最模范的妇女，以便在那里过夜；我们请她给我们一个完全属于自己的地方，没有妇女或童贞女进入。这位老妇人亲自给我们带来灯，并亲自带来我们所需的一切。出于对弟兄们的爱，她会带来服侍陌生弟兄所需的一切。当她睡觉的时刻来到时，她会平静地离开，回她的家去。\par

% structure: p0039 source=h3 target=subsection reason=relative-heading-rank; base=h2

\subsection{第五章 在只有一位妇女之处，教父不逗留；应如何谨慎避免绊脚石}

但是，倘若我们偶然来到一个地方，那里只有一位信主的女子，别无他人，我们便不停留，也不祈祷，也不宣读圣经，而是如逃避蛇的面容、逃避罪的面容一般逃离。并非我们藐视那信主的女子——决不至于我们对基督内的弟兄怀有这样的心思！——而是因为她单独一人，我们害怕有人会用虚假的言语对我们进行诽谤。因为人的心坚定于恶。并且，为使我们不给那些寻找借口攻击我们、毁谤我们的人留下借口，也不使我们成为任何人的绊脚石，因此我们断绝了那些寻找借口之人的借口；因此我们必须\BibleQuote{提防，不使自己成为任何人的绊脚石，无论是犹太人、外邦人，还是天主的教会；我们不应只求自己的益处，而应求众人的益处，使他们得救。}\BibleRef{格前 10:32}因为别人因我们而跌倒，对我们并无益处。所以，我们应当时刻谨慎提防，不使我们的弟兄受伤，不因我们成为他们的绊脚石而使他们喝下有扰良心之物。因为\BibleQuote{若因食物使我们的弟兄忧愁、震惊、软弱或跌倒，我们便不在天主的爱中行走。你因食物使基督为他而死的人丧亡。}因为，\Emph{在}\BibleQuote{这样得罪你们的弟兄，击伤他们软弱的良心，你们就是得罪基督自己。因为，若因食物我的弟兄跌倒，}我们\Emph{身为信徒的}要说：\BibleQuote{我们永远不再吃肉，免得使我们的弟兄跌倒。}\BibleRef{格前 8:12}这些事，凡是真正爱天主、真正背起自己的十字架、穿上基督并爱邻人的人，都会这样做；那人谨慎自己，不成为任何人的绊脚石，不使人因他的缘故跌倒，不因他常与少女同住同屋——这是不当的——以致看见和听见的人覆亡。这样的恶行充满绊跌和危险，是与死亡相近的。但那为了纯洁而在一切事上谨慎畏惧的人是有福的！\par

% structure: p0041 source=h3 target=subsection reason=relative-heading-rank; base=h2

\subsection{第六章 基督徒在外邦人中应当如何行事。}

此外，若我们偶然去一个没有基督徒的地方，并且必须在那里停留数日，我们应\BibleQuote{机警如同蛇，纯朴如同鸽子；}\BibleRef{玛 10:16}并应\BibleQuote{不要像愚笨人，但要像明智人，}\BibleRef{弗 5:15}怀着对天主敬畏的一切\Emph{自制}，好使天主因我们贞洁圣洁的行为，藉着我们的主耶稣基督，在一切事上受光荣。因为，\BibleQuote{你们或吃或喝，或无论做什么，一切都要为光荣天主而做。}\BibleRef{格前 10:31}愿\BibleQuote{凡看见我们的，都承认我们是蒙祝福的苗裔，}\BibleRef{依 61:9}\BibleQuote{永生天主的儿女，}\BibleRef{斐 2:15}在一切事上——在\Emph{我们的}所有话语中，存有羞怯、纯洁、谦卑，因我们不效法外邦人任何事，也不\Emph{作为信徒}像\Emph{其他}人那样，而是在一切事上与恶者疏远。我们\BibleQuote{不把圣物投给狗，也不把珍珠丢给猪；}\BibleRef{玛 7:6}却以一切可能的\Emph{自制}、一切谨慎、一切对天主的敬畏和热切的心神赞美天主。因为，我们不在外邦人醉酒、以污秽言语在他们宴席上亵渎的地方事奉，因他们的邪恶。因此，我们不为外邦人歌唱\Emph{圣咏}，也不向他们诵读圣经，免得我们像\Emph{普通的}歌者，无论是弹竖琴的，或是用声音歌唱的，或是\Emph{像}占卜者，有许多人这样做，为了填饱可怜的几口面包，为了一杯劣酒，\Emph{到处}去\BibleQuote{在外邦人的异乡歌唱主的歌曲}，做不当之事。我弟兄们，不要这样；我们恳求你们，我弟兄们，不要让这些事在你们中间发生；但要移除那些选择以如此无耻和可耻方式行事的人。我弟兄们，这些事不正当。但我们恳求你们，正义中的弟兄们，愿这些事在你们中间如同在我们中间一样\Emph{实行}，作为信徒和将要相信之人的模范。让我们属于基督的羊群，在一切正义、一切圣洁无瑕的品行中，行为正直圣善，正如信徒所应当的，并遵守那些可赞美、纯洁、圣洁、可敬、高尚的事；你们要促进一切有益的事。因为你们是\BibleQuote{我们的喜乐、我们的冠冕，}我们的希望，我们的生命，\BibleQuote{只要你们在主内站稳。}\BibleRef{斐 4:1}但愿如此！\par

% structure: p0043 source=h3 target=subsection reason=relative-heading-rank; base=h2

\subsection{第七章 思考警戒的榜样与教导的模式之益处。}

因此，我弟兄们，让我们思考并察看所有正义的父辈们在\Emph{此}生寄居期间是如何行事的，让我们从法律\Emph{直到}新约进行探索和查考。因为这是既合宜又有益的，我们应知道有多少人，以及\Emph{他们是谁}，因妇女而丧亡；又有多少妇女，以及\Emph{她们是谁}，因男人而丧亡，由于他们彼此经常交往。并且，同样地，我还要指出有多少人，以及\Emph{他们是谁}，终身彼此共处，甚至到死，一直\Emph{行}无瑕的贞洁之工。这是明显且众所周知的。\par

% structure: p0045 source=h3 target=subsection reason=relative-heading-rank; base=h2

\subsection{第八章 若瑟与普提法尔之妻；对女性的爱应当如何。}

\Emph{有}若瑟此人，忠实、聪慧、明智，凡事敬畏天主。岂不是一个妇人对这贞洁正直之人的美貌产生了过度的情欲？当他拒绝并不同意满足她的强烈欲望时，她便\Emph{作}假见证，将这义人投入各种困苦折磨，几乎致死。但天主救他脱离这\Emph{个}可怜妇人所带来的一切祸患。弟兄们，你们看，时常看见那埃及妇人的面容给这义人带来了何等痛苦。因此，我们不可常与妇女或少女相处。这对那些真心愿意\BibleQuote{束上腰}的人是无益的。\BibleRef{路 12:35} 因为要求我们以完全的纯洁和贞洁，并完全约束思想，在敬畏天主中爱姊妹，不可常与她们往来，也不可时刻接近她们。\par

% structure: p0047 source=h3 target=subsection reason=relative-heading-rank; base=h2

\subsection{第九章 三松的警戒性堕落。}

你难道没有听说过纳齐尔人三松吗？\BibleQuote{天主的灵与他同在。} \BibleRef{民 13:25} 这个大有能力的人？此人本为纳齐尔人，奉献于天主，并\Emph{赋有}力量和强健，一个妇人却用\Emph{她的}可怜身体和\Emph{她的}卑劣情欲将他毁掉。难道你或许与他一样吗？认识你自己，认识你力量的程度。\BibleQuote{有夫之妇猎取宝贵的生命。} \BibleRef{箴 6:26} 因此，我们不允许任何人与已婚妇人同坐，更不用说与立誓守贞的少女同住一屋，或睡在她睡的地方，或常与她一起。因为敬畏天主之人应当憎恶厌弃这等事。\par

% structure: p0049 source=h3 target=subsection reason=relative-heading-rank; base=h2

\subsection{第十章 达味的罪，对我们软弱之人的警戒。}

达味的例子难道不教训你吗？天主\Quote{找到他合祂心意的人}，一个忠实、无过、虔诚、真实的人？他看见了一个妇人的美貌——我指的是巴特舍巴——看见她正在无人时洗澡净身。这圣者见了这妇人，便完全被其美貌所俘虏。请看，因着一个妇人，他犯下何等恶行，这位义人\Emph{如何}犯罪，并下令在战场上杀死这妇人的丈夫。你看见了他谋划执行了何等邪恶的计谋，以及他\Emph{如何}因对妇人的情欲而犯下谋杀——\Emph{他}，达味，被称为\Quote{主的受傅者}。人啊，你要受劝诫：因为，如果这样的人都被妇人毁掉，你的\Emph{义德}算什么？你在圣人中算什么？你竟毫无敬畏地昼夜与妇女少女愚昧相处。不\Emph{如此}，我的弟兄们，不如此，而要记住有关妇人的话：\BibleQuote{她的手是罗网，她的心是陷阱；义人将逃脱她，恶人却落入她手中。} \BibleRef{训 7:26} 因此，我们这些奉献者要小心，不可与立誓守贞的妇女同住一屋。因为这样的行为对天主的仆人不合适，也不正当。\par

% structure: p0051 source=h3 target=subsection reason=relative-heading-rank; base=h2

\subsection{第十一章 达味乱伦子女的警戒史。}

你难道没有读过达味的子女阿默农和塔玛尔的事吗？这阿默农对他的妹妹起了情欲，并玷辱了她，没有怜惜她，因为他以可耻的欲火恋慕她；由于常与她来往，没有敬畏天主，他成了邪恶放荡之人，\BibleQuote{在以色列中行了丑事。} \BibleRef{创 34:7} 因此，我们与姊妹来往，\Emph{沉溺}于欢笑和放荡，是不合适也不应当的；但\Emph{我们对待她们}应当极其贞洁纯净，且敬畏主。\par

% structure: p0053 source=h3 target=subsection reason=relative-heading-rank; base=h2

\subsection{第十二章 撒罗满因妇人而痴迷。}

你难道没有读过撒罗满的历史吗？他是达味之子，天主赐给他智慧、见识、豁达的心、财富和极大的荣耀，超越众人。然而这同一个人，因着妇人，走向了毁灭，离弃了主。\par

% structure: p0055 source=h3 target=subsection reason=relative-heading-rank; base=h2

\subsection{第十三章 苏撒纳的史实教训我们眼目与交际要谨慎。}

你难道没有读过，也不知道苏撒纳时代的那两位长老吗？他们因为常与妇人在一起，注视他人的美貌，陷入极度的淫荡，不能保守纯洁的心，反被邪僻的性情所胜，突然闯到有福的苏撒纳面前要玷污她。但她没有同意他们丑恶的情欲，反而呼求天主，天主就救她脱离了恶长老之手。所以，难道我们不该战栗惧怕吗？因为这些老人，本是天主子民的审判官和长老，却因一个妇人从高位坠落。因为他们没有记住所说的：\Quote{不可贪恋他人的美貌；} 以及\BibleQuote{妇人的美貌毁灭了许多人；} \BibleRef{德 9:8} 和\BibleQuote{不可与已婚妇人同坐；} \BibleRef{德 9:12} 以及\Emph{经上}又说：\BibleQuote{有人怀里藏火，衣服能不烧着吗？} \BibleRef{箴 6:27} 或\BibleQuote{人若在火上行走，脚岂能不被烫伤？所以，凡进入他人妻子的，不能无罪；凡亲近她的，必不能逃脱。} \BibleRef{箴 6:28} 又经上说：\BibleQuote{不可贪恋妇人的美貌，免得你被她的眼睑捕获；} \BibleRef{箴 6:25} 以及\BibleQuote{不可注视少女，免得你因爱慕她而丧亡；} \BibleRef{德 9:5} 和\BibleQuote{不可常与歌女相处；} \BibleRef{德 9:4} 并\BibleQuote{自以为站得稳的，要小心，免得跌倒。} \BibleRef{格前 10:12}\par

% structure: p0057 source=h3 target=subsection reason=relative-heading-rank; base=h2

\subsection{第十四章 旧约中谨慎行为的例子。}

但请看经上也论及那些圣人，即先知们，以及我们主的宗徒们。让我们看看这些圣人中，是否有人常与少女、或年轻的已婚妇女、或神圣的宗徒拒绝接纳的寡妇在一起。让我们怀着敬畏天主的心，思量这些圣人的生活模式。看！我们发现经上记载\Person{梅瑟}和\Person{亚郎}，他们行事和生活都与那些同样遵循相同行为方式的男人为伴。同样，\Person{农}的儿子\Person{若苏厄}也是如此。没有女子与他们同在；他们独自圣洁地在天主面前服务，男人与男人在一起。不仅如此，他们还教导百姓，每当营队移动时，每个支派应分开移动，妇女与妇女分开，且应走在营队后面，男人也按支派分开。他们遵照主的命令出发，如同一个智慧的民族，使营队移动时不因妇女而混乱。他们以美丽有序的排列前进，绝不绊跌。因为看！圣经为我的话作证：\BibleQuote{当以色列子民渡过\Place{苏特海}后，\Person{梅瑟}和以色列子民歌颂主说：我们要赞美主，因为祂极其可颂。}\newline{}\BibleRef{出 15:1} \Person{梅瑟}唱完赞歌之后，\Person{梅瑟}和\Person{亚郎}的姐妹\Person{米黎盎}手拿手鼓，所有妇女都跟着她出来，与她一起歌唱赞美，妇女与妇女分开，男人与男人分开。再者，我们发现\Person{厄里叟}、\Person{革哈齐}和众先知之子一同生活在敬畏主之中，没有女子与他们同住。\Person{米该亚}也\Emph{如此}，所有先知也同样生活在敬畏主之中。\par

% structure: p0059 source=h3 target=subsection reason=relative-heading-rank; base=h2

\subsection{第15章 耶稣的榜样；我们如何可以允许自己被妇女服侍}

而且，为免我们的论述\Emph{过于}冗长，关于我们的主\Person{耶稣基督}，我们能说什么呢？我们的主亲自与祂的十二位门徒常在一起，当祂来到世间时。不仅如此，当祂派遣他们时，祂让他们两个两个一起出发，男人与男人；但妇女没有被派遣与他们同去，无论在路边或在家里，他们都不与妇女或少女交往：如此他们在一切事上中悦天主。此外，当我们的主\Person{耶稣基督}独自在井旁与\Place{撒玛黎雅}妇人谈话时，\BibleQuote{祂的门徒来了}，见祂正同那妇人谈话，\BibleQuote{就惊奇\Person{耶稣}正站着与一个妇人谈话。}\newline{}\BibleRef{若 4:27} 难道这不是一个不可更改的规则，一个榜样，一个模式，给所有人类族群吗？不仅如此，当我们的主从死人中复活后，\Person{玛利亚}来到坟墓处，她跑去抱住我们主的脚，朝拜祂，并想抓住祂。但祂对她说：\BibleQuote{你别拉住我不放，因为我还没有升到父那里去。}\newline{}\BibleRef{若 20:17} 那么，这不是令人惊讶吗？我们的主不允许那位有福的妇女\Person{玛利亚}触摸祂的脚，而\Emph{你}却与她们同住，由妇女和少女服侍，睡在她们睡觉的地方，妇女为你洗脚，给你傅油！啊，这种该受责备的心态！啊，这种缺乏敬畏的心态！啊，这种厚颜和愚昧，没有敬畏天主！你难道不判断自己吗？你难道不省察自己吗？你难道不知道你自己和你力量的尺度吗？此外，这些事是可信的，这些事是真实和正确的；这些是对那些在我们主内正直行事之人不可改变的规则。还有许多圣洁的妇女曾用她们的财物服侍圣洁的人，就如叔能妇人服侍\Person{厄里叟}；但她没有与他同住，先知单独住在一所房子里。当她的儿子死后，她想要仆倒在先知脚前；但他的仆人不允许她，而是阻止了她。但\Person{厄里叟}对仆人说：\BibleQuote{让她去吧，因为她心里很痛苦。}\newline{}\BibleRef{列下 4:27} 因此，从这些事我们应当理解他们的生活方式。妇女们曾用她们的财物服侍我们的主\Person{耶稣基督}；但她们没有与祂同住；而是纯洁、圣洁、无可指责地生活在主面前，完成她们的赛程，在我们的主全能天主那里接受了冠冕。\par

% structure: p0061 source=h3 target=subsection reason=relative-heading-rank; base=h2

\subsection{第16章 劝勉合一与服从；结论}

因此，我们恳求你们，我们在主内的弟兄们，愿这些事在你们中间也如在我们中间一样被遵守，使我们同心合意，使我们在你们内，你们在我们内成为一体，并在一切事上，我们能在主内同心同德。凡认识主的，必听从我们；凡不属于天主的，必不听从我们。那真正渴望持守圣洁的，必听从我们；那真正渴望持守童贞的童贞女，必听从我们；但那不真正渴望持守童贞的，必不听从我们。最后，在主内安好，你们所有圣人，要在主内喜乐。愿平安与喜乐从天主父经由我们的主\Person{耶稣基督}与你们同在。阿们。\par

\Emph{这里}结束\Person{克来孟}——\Person{伯多禄}的门徒的第二封书信。愿他的祈祷与我们同在！阿们。\par

\SourceRecord{Translated by B.P. Pratten.；Ante-Nicene Fathers；Vol. 8.；Edited by Alexander Roberts, James Donaldson, and A. Cleveland Coxe.；Buffalo, NY: Christian Literature Publishing Co.,；1886.；<http://www.newadvent.org/fathers/0803.htm>.}
